\section{Bézier Curves}

\begin{frame}{Quadratic Bézier: LERP of LERPs}
  \begin{columns}
    \begin{column}{0.8\textwidth}
      \begin{mathbox}{Construction Method}
        \small
        Given three control points $\mathbf{P_0}$, $\mathbf{P_1}$, $\mathbf{P_2}$:

        \vspace{0.3cm}
        \textbf{Step 1:} Linear interpolation
        \begin{align*}
          \mathbf{Q_0}(t) &= \text{lerp}(\mathbf{P_0}, \mathbf{P_1}, t) \\
          \mathbf{Q_1}(t) &= \text{lerp}(\mathbf{P_1}, \mathbf{P_2}, t)
        \end{align*}

        \textbf{Step 2:} Another linear interpolation
        \begin{align*}
          \mathbf{B}(t) &= \text{lerp}(\mathbf{Q_0}(t), \mathbf{Q_1}(t), t) \\
          &= (1-t)^2\mathbf{P_0} + 2t(1-t)\mathbf{P_1} + t^2\mathbf{P_2}
        \end{align*}

        \textbf{Definition:}
        \[
          \text{quadBezier}(\mathbf{A},\mathbf{B},\mathbf{C},t)
          = (1-t)^2\mathbf{A} + 2t(1-t)\mathbf{B} + t^2\mathbf{C}
        \]
      \end{mathbox}
    \end{column}

    \begin{column}{0.3\textwidth}
      \begin{center}
        \begin{tikzpicture}[scale=0.8]
          % Control points
          \coordinate (P0) at (0.5,0.5);
          \coordinate (P1) at (2,2.5);
          \coordinate (P2) at (3.5,1);

          % Control polygon
          \draw[dashed, LightGray] (P0) -- (P1) -- (P2);

          % Bézier curve
          \draw[thick, PrimaryColor] (P0) .. controls (P1) .. (P2);

          % Control points
          \fill[SecondaryColor] (P0) circle (3pt);
          \fill[SecondaryColor] (P1) circle (3pt);
          \fill[SecondaryColor] (P2) circle (3pt);

          \node[below left] at (P0) {\scriptsize $\mathbf{P_0}$};
          \node[above] at (P1) {\scriptsize $\mathbf{P_1}$};
          \node[below right] at (P2) {\scriptsize $\mathbf{P_2}$};

          % Example t interpolation
          \coordinate (Q0) at (1.25,1.5);
          \coordinate (Q1) at (2.75,1.75);
          \coordinate (B) at (2,1.625);

          % Intermediate points
          \fill[AccentColor] (Q0) circle (2pt);
          \fill[AccentColor] (Q1) circle (2pt);
          \fill[ObjectColor] (B) circle (3pt);

          % Intermediate lines
          \draw[AccentColor] (Q0) -- (Q1);

          \node[left] at (Q0) {\scriptsize $\mathbf{Q_0}$};
          \node[right] at (Q1) {\scriptsize $\mathbf{Q_1}$};
          \node[below] at (B) {\scriptsize $\mathbf{B}(t)$};
        \end{tikzpicture}
      \end{center}
    \end{column}
  \end{columns}

  \vspace{0.1cm}
  \begin{center}
    \textcolor{AccentColor}{\href{https://www.desmos.com/calculator/j0cmdymljy}{[Desmos Demo: Quadratic Bézier Construction]}}
  \end{center}
\end{frame}

\begin{frame}{Quadratic Bézier: Mathematical Form}
  \begin{mathbox}{Expanded Form}
    \small
    Starting from the LERP of LERPs:
    \begin{align*}
      \mathbf{B}(t) &= (1-t)[(1-t)\mathbf{P_0} + t\mathbf{P_1}] + t[(1-t)\mathbf{P_1} + t\mathbf{P_2}] \\
      &= (1-t)^2\mathbf{P_0} + 2t(1-t)\mathbf{P_1} + t^2\mathbf{P_2}
    \end{align*}

    \textbf{Bernstein Form:}
    \begin{align*}
      \mathbf{B}(t) &= B_{0,2}(t)\mathbf{P_0} + B_{1,2}(t)\mathbf{P_1} + B_{2,2}(t)\mathbf{P_2}
    \end{align*}

    Where the Bernstein polynomials are:
    \begin{align*}
      B_{0,2}(t) &= (1-t)^2 & B_{1,2}(t) &= 2t(1-t) & B_{2,2}(t) &= t^2
    \end{align*}
  \end{mathbox}

  \begin{conceptbox}{Key Properties}
    \footnotesize
    \begin{itemize}
      \item \textbf{Endpoint Interpolation:} $\mathbf{B}(0) = \mathbf{P_0}$, $\mathbf{B}(1) = \mathbf{P_2}$
      \item \textbf{Tangent Control:} $\mathbf{P_1}$ controls the tangent direction
        % \item \textbf{Convex Hull:} Curve lies within triangle formed by control points
    \end{itemize}
  \end{conceptbox}
\end{frame}

\begin{frame}{Cubic Bézier: LERP of Quadratics}
  \begin{columns}
    \begin{column}{0.65\textwidth}
      \begin{mathbox}{Construction Method}
        \small
        Four control points $\mathbf{P_0}$, $\mathbf{P_1}$, $\mathbf{P_2}$, $\mathbf{P_3}$:

        \vspace{0.3cm}
        \textbf{Step 1:} Build two quadratic Béziers
        \begin{align*}
          \mathbf{Q_0}(t) &= \text{quadBezier}(\mathbf{P_0}, \mathbf{P_1}, \mathbf{P_2}, t) \\
          \mathbf{Q_1}(t) &= \text{quadBezier}(\mathbf{P_1}, \mathbf{P_2}, \mathbf{P_3}, t)
        \end{align*}

        \textbf{Step 2:} Interpolate between them
        \begin{align*}
          \mathbf{B}(t) &= \text{lerp}(\mathbf{Q_0}(t), \mathbf{Q_1}(t), t) \\
          &= (1-t)^3\mathbf{P_0} + 3t(1-t)^2\mathbf{P_1} \\
          &\quad+ 3t^2(1-t)\mathbf{P_2} + t^3\mathbf{P_3}
        \end{align*}
      \end{mathbox}
    \end{column}

    \begin{column}{0.35\textwidth}
      \begin{center}
        \begin{tikzpicture}[scale=0.8]
          % Control points
          \coordinate (P0) at (0.5,1);
          \coordinate (P1) at (1,2.5);
          \coordinate (P2) at (3,0.5);
          \coordinate (P3) at (3.5,2);

          % Control polygon
          \draw[dashed, LightGray] (P0) -- (P1) -- (P2) -- (P3);

          % Cubic Bézier curve
          \draw[thick, PrimaryColor] (P0) .. controls (P1) and (P2) .. (P3);

          % Control points
          \fill[SecondaryColor] (P0) circle (3pt);
          \fill[AccentColor] (P1) circle (3pt);
          \fill[AccentColor] (P2) circle (3pt);
          \fill[SecondaryColor] (P3) circle (3pt);

          \node[below] at (P0) {\scriptsize $\mathbf{P_0}$};
          \node[above] at (P1) {\scriptsize $\mathbf{P_1}$};
          \node[below] at (P2) {\scriptsize $\mathbf{P_2}$};
          \node[above] at (P3) {\scriptsize $\mathbf{P_3}$};

          % Example t
          \coordinate (Q0) at (1.25,1.75); % point on first quadratic
          \coordinate (Q1) at (2.75,1.25); % point on second quadratic
          \coordinate (B) at (2,1.5);      % point on cubic

          % Quadratic points
          \fill[AccentColor] (Q0) circle (2pt);
          \fill[AccentColor] (Q1) circle (2pt);
          \fill[ObjectColor] (B) circle (3pt);

          \draw[AccentColor] (Q0) -- (Q1);

          \node[left] at (Q0) {\scriptsize $\mathbf{Q_0}$};
          \node[right] at (Q1) {\scriptsize $\mathbf{Q_1}$};
          \node[above] at (B) {\scriptsize $\mathbf{B}(t)$};
        \end{tikzpicture}
      \end{center}
    \end{column}
  \end{columns}

  \vspace{0.5cm}
  \begin{center}
    \textcolor{AccentColor}{\href{https://www.desmos.com/calculator/hlsbm9h3bo}{[Desmos Demo: Cubic Bézier as LERP of Quadratics]}}
  \end{center}
\end{frame}
